% =====================================================================
%  4.6  Atividade 3.2.2.6 -- Programa Sanepar Rural
% =====================================================================
\subsection{Atividade 3.2.2.6 -- Implantar sistemas coletivos de abastecimento de água em comunidades rurais atendidas pela Sanepar (Programa Sanepar Rural)}
\label{subsec:sanepar-rural}

\subsubsection{Objetivo e fundamentação}

Implementar o abastecimento de água potável em comunidades rurais dos
municípios que têm Contrato de Concessão ou de Programa vigente com a Sanepar,
segundo os critérios do Programa Sanepar Rural. Para isso, a Sanepar
fornecerá expertise, estudos técnicos, apoio técnico, ambiental e
sociocomunitário, treinamento, materiais hidráulicos e equipamentos
\cite{mop2026}.

\begin{fichaatividade}{Ficha-síntese da Atividade 3.2.2.6}{qua:ficha-3226}
  \campo{Código}{3.2.2.6}
  \campo{Tipo de atividade}{Serviço e obra}
  \campo{Forma de execução}{Convênio}
  \campo{Fonte de recursos}{Contrapartida da Sanepar}
  \campo{Responsável}{Sanepar / IAT}
  \campo{Dependências institucionais}{Prefeituras}
  \campo{Prazo estimado}{Ano 1 ao ano 5}
  \campo{Produto esperado}{Sistemas implantados}
  \campo{Unidade de medida}{Número de pessoas atendidas}
  \campo{Evidência de comprovação}{Relatório técnico aprovado}
  \campo{Periodicidade}{Contínua}
  \campo{Validação}{Sanepar}
\end{fichaatividade}

\subsubsection{Procedimentos}

As atividades seguirão o manual do Programa Sanepar Rural \cite{saneparrural},
com destaque para:

\begin{enumerate}[label=\alph*)]
  \item levantamento técnico e estudo técnico de engenharia;
  \item formalização da parceria por meio de Termo de Compromisso e
        Responsabilidade ao Contrato de Concessão ou de Programa;
  \item elaboração dos projetos de engenharia dos sistemas de abastecimento
        de água;
  \item fornecimento dos materiais hidráulicos necessários à obra;
  \item apoio técnico ao município na execução das obras;
  \item mobilização e organização comunitária, com apoio à formação de
        associações de moradores;
  \item treinamento e educação ambiental: uso racional da água, higiene e
        manutenção preventiva dos sistemas;
  \item monitoramento e suporte técnico pós-obra, para garantir a qualidade
        da água e o suporte operacional às associações.
\end{enumerate}

\pendencia{A meta desta atividade no Quadro 31 do MOP (59.780 pessoas) é a
meta total do indicador ``Pessoas atendidas com água gerida de forma segura''
do Componente~3, que soma os sistemas do IAT (cerca de 16.200 pessoas) e da
Sanepar (cerca de 43.500). A parcela da Sanepar Rural é de cerca de 43.500
pessoas. Além disso, o MOP indica periodicidade ``anual'' no texto e
``contínua'' no Quadro 31.}

\pendencia{Falta um quadro de etapas com prazos, produtos e evidências, como
nas demais atividades. O manual do Programa Sanepar Rural deve ser citado com
título, versão e data.}

\subsubsection{Interfaces}

\begin{itemize}
  \item \atividade{3.2.2.1}: os modelos de gestão também se aplicam aos
        sistemas implantados pela Sanepar Rural.
  \item Indicador de resultado do Componente~3 ``Pessoas atendidas com água
        gerida de forma segura'' (\autoref{tab:metas-componente3}).
\end{itemize}
